AssistVision turns a pair of ordinary spectacles into a meal-intake monitor for
older adults: a clip-on camera records the plate when the meal is served and
again when it is cleared, the host works out the share eaten, family members
are alerted after a run of low-intake meals, and the doctor gets a weekly
report. The parts cost about INR~2,000--2,900 and two of them are printed. The
vision experiments explain the design. A single image anchored only on the
table does not recover food volume: the best deployable models (ZoeDepth-NYU,
and Depth Pro at 15$\times$ the latency) still have median errors of
44--48\,\%, and low depth error (AbsRel) does not imply low volume error.
Several views with marker scale work on textured synthetic food on a
turntable path (E5). From the glasses' short head sweep (E6) the share eaten is off
by a median of 8 points on steel and 16.5 on white ceramic, against 2.8 on the
turntable path, and only 5 of 12 pairs come within 10 points. Textureless food fails in both cases, and
volume does not predict grams on real mixed dishes, which is why the glasses
report a share rather than grams. The same two photos also show which dish
was left: on 6 simulated multi-dish plates the dish left most was named
correctly each time, and the app turns that into a protein-gap alert that says
whether the dal is being left or the menu is short. Next we will print and wear
the glasses, capture real thalis against weighed plate waste, name the dishes
automatically, and pilot the app with residents, families and staff in a care
home.
